\chapter{Der Leitfall, zu Ende gerechnet}

\section{Alle Schätzungen nebeneinander}

\begin{merksatz}
\boxtitel{Zehn Antworten auf eine Frage}
Wie stark senkt ein Euro Zuzahlung die Zahl der Arztbesuche je Jahr?

\begin{center}
{\small
\renewcommand{\arraystretch}{1.46}
\begin{tabular}{@{}P{0.8cm}P{3.2cm}P{2.0cm}P{2.6cm}@{}}
\toprule
\textbf{Kap.} & \textbf{Verfahren} & \textbf{Wert} & \textbf{trägt auf}\\
\midrule
$1$ & kleinste Quadrate & $-0{,}376$ & nichts\\
$1$ & mit Gesundheit & $-0{,}098$ & Messung der Störgröße\\
$1$ & ausgelost & $-0{,}100$ & Randomisierung\\
$3$ & Instrument, stark & $-0{,}100$ & Ausschluss\\
$3$ & Instrument, schwach & $-0{,}8$ bis $0{,}9$ & ---\\
$4$ & GMM, zwei Instr. & $-0{,}098$ & Ausschluss, Test\\
$7$ & feste Effekte & $-0{,}099$ & Konstanz der Störgröße\\
$9$ & doppelte Differenz & $-0{,}100$ & parallele Trends\\
$9$ & gestaffelt, zweiseitig & zu klein & \textcolor{akzent}{verbotene Vergleiche}\\
$13$ & unter neuer Regel & $-0{,}141$ & Lucas\\
\bottomrule
\end{tabular}}
\end{center}

\begin{center}
\begin{minipage}{0.9\textwidth}
\textbf{Sechs Verfahren treffen den wahren Wert --- jedes mit einer anderen Annahme. Keines ist ohne.}
\end{minipage}
\end{center}

\begin{center}
\begin{minipage}{0.88\textwidth}
\kursiv{\textbf{Die rechte Spalte ist die eigentliche Tabelle.} Dass die sechs übereinstimmen, ist in der Simulation selbstverständlich, weil alle Annahmen dort gelten. In Wirklichkeit stimmen sie nicht überein, und dann sagt die rechte Spalte, worüber man streitet: Ist die Störgröße konstant? Sind die Trends parallel? Wirkt das Instrument nur über die Zuzahlung? Die Ökonometrie hat keine Methode, die diese Fragen beantwortet --- sie hat Methoden, die sie stellen.}
\end{minipage}
\end{center}
\end{merksatz}

\section{Was das Experiment fand}

\begin{merksatz}
\boxtitel{Das Krankenversicherungsexperiment, in Zahlen}
$1974$ bis $1982$, rund $7700$ Teilnehmer, zufällig Tarifen mit $0$, $25$, $50$ oder $95$ Prozent Selbstbeteiligung zugewiesen, mit einer Obergrenze je Haushalt.

\begin{center}
{\small
\renewcommand{\arraystretch}{1.48}
\begin{tabular}{@{}P{3.0cm}P{2.2cm}P{2.2cm}@{}}
\toprule
\textbf{Selbstbeteiligung} & \textbf{Ausgaben} & \textbf{relativ}\\
\midrule
$0\,\%$ & $100$ & ---\\
$25\,\%$ & $81$ & $-19\,\%$\\
$50\,\%$ & $77$ & $-23\,\%$\\
$95\,\%$ & $\mathbf{69}$ & $\mathbf{-31\,\%}$\\
\bottomrule
\end{tabular}}
\end{center}
\kursiv{Gerundete Verhältnisse nach den veröffentlichten Ergebnissen; die Preiselastizität lag bei etwa $-0{,}2$.}

\schritt{Vier Befunde}
\begin{center}
{\small
\renewcommand{\arraystretch}{1.5}
\begin{tabular}{@{}P{1.0cm}P{7.4cm}@{}}
\toprule
$1$ & Zuzahlung senkt die Inanspruchnahme deutlich --- vor allem die Zahl der Behandlungsepisoden, kaum die Kosten je Episode.\\
$2$ & Sie senkt nötige und unnötige Leistungen gleichermaßen --- Patienten unterscheiden nicht.\\
$3$ & Die Gesundheit der meisten Teilnehmer litt nicht messbar.\\
$4$ & Bei armen Teilnehmern mit Bluthochdruck, Sehschwäche und Zahnproblemen litt sie.\\
\bottomrule
\end{tabular}}
\end{center}

\begin{center}
\begin{minipage}{0.88\textwidth}
\kursiv{\textbf{Der vierte Befund ist der, der im politischen Streit meist fehlt.} Er ist ein heterogener Effekt im Sinne von Band 23, Kapitel~14 --- und das Experiment konnte ihn nur finden, weil es groß genug war und die Gesundheit tatsächlich maß. Die Elastizität von $-0{,}2$ wird bis heute in jeder Rechnung zu Zuzahlungen benutzt; sie gilt für amerikanische Familien der 1970er Jahre mit einer Obergrenze der Jahresbelastung, und das ist die Übertragbarkeitsfrage aus Band 23 in ihrer schwierigsten Form.}
\end{minipage}
\end{center}
\end{merksatz}

\section{Was die Beobachtungsstudien fanden}

\begin{merksatz}
\boxtitel{Die Praxisgebühr}
Deutschland, $2004$ bis $2012$: zehn Euro je Quartal beim ersten Arztbesuch, befreit waren Kinder und Versicherte über einer Belastungsgrenze.

\begin{center}
{\small
\renewcommand{\arraystretch}{1.5}
\begin{tabular}{@{}P{3.0cm}P{2.4cm}P{3.0cm}@{}}
\toprule
\textbf{Entwurf} & \textbf{Befund} & \textbf{Annahme}\\
\midrule
vorher/nachher & $-8$ bis $-10\,\%$ & kein Trend\\
doppelte Differenz, Kinder & etwa $-7\,\%$ & Kinder parallel\\
doppelte Differenz, Befreite & kaum Effekt & Befreite parallel\\
Schwelle an Belastungsgrenze & kleiner Effekt & keine Sortierung\\
\bottomrule
\end{tabular}}
\end{center}

\begin{center}
\begin{minipage}{0.9\textwidth}
\textbf{Die Befunde streuen von null bis zehn Prozent, und die Streuung geht auf die Wahl der Kontrollgruppe zurück --- nicht auf die Daten.}
\end{minipage}
\end{center}

\begin{center}
\begin{minipage}{0.88\textwidth}
\kursiv{\textbf{Jede Kontrollgruppe hat ihren Fehler}: Kinder wurden in denselben Jahren häufiger zu Vorsorgeuntersuchungen geschickt; Befreite sind chronisch krank und reagieren anders; an der Belastungsgrenze sortieren sich Versicherte, die ihre Ausgaben kennen. \textbf{Das ist die Tabelle aus Abschnitt 14.1 mit wirklichen Daten}: Die rechte Spalte entscheidet, und sie ist nirgends unstrittig. Als die Gebühr 2012 abgeschafft wurde, stiegen die Besuche nicht messbar --- was entweder heißt, dass sie nie gewirkt hatte, oder dass die Gewöhnung blieb. Auch das ist eine Frage der Kontrollgruppe.}
\end{minipage}
\end{center}
\end{merksatz}

\section{Wie man es berichtet}

\begin{merksatz}
\boxtitel{Die Berichtspflicht}
\begin{center}
{\small
\renewcommand{\arraystretch}{1.52}
\begin{tabular}{@{}P{2.8cm}P{5.8cm}@{}}
\toprule
\textbf{Die Frage} & Welcher Effekt, für wen, unter welcher Politik --- Band 23, Kapitel~2, und Kapitel~13 dieses Bandes.\\
\textbf{Die Annahme} & Die eine, auf der alles ruht --- in einem Satz, nicht in einer Fußnote.\\
\textbf{Die Prüfung} & Was man geprüft hat (Vorperioden, erste Stufe, Überidentifikation) und was nicht prüfbar war.\\
\textbf{Die Empfindlichkeit} & Wie weit die Annahme verletzt sein dürfte, bevor das Vorzeichen kippt.\\
\textbf{Die Spanne} & Mehrere Verfahren nebeneinander, mit ihren Annahmen --- nicht das günstigste.\\
\bottomrule
\end{tabular}}
\end{center}

\begin{center}
\begin{minipage}{0.9\textwidth}
\textbf{Eine Schätzung ohne ihre Annahme ist eine Zahl. Eine Schätzung mit ihrer Annahme ist ein Befund.}
\end{minipage}
\end{center}

\begin{center}
\begin{minipage}{0.88\textwidth}
\kursiv{\textbf{Die fünfte Zeile ist die wichtigste und die am seltensten befolgte.} Wer zehn Spezifikationen gerechnet und eine berichtet hat, betreibt die Suche aus Band 17, Kapitel~12 --- nur mit Modellen statt mit Hypothesen. Die Ökonometrie hat dafür seit einigen Jahren das Werkzeug der Spezifikationskurve: alle vertretbaren Varianten nebeneinander, geordnet nach Effektgröße, und man sieht, ob das Vorzeichen hält.}
\end{minipage}
\end{center}
\end{merksatz}

\section{Der Bruchfall}

\begin{warnung}
\boxtitelr{Was alle Verfahren teilen}
\begin{center}
{\small
\renewcommand{\arraystretch}{1.52}
\begin{tabular}{@{}P{2.8cm}P{5.8cm}@{}}
\toprule
\textbf{Linearität} & Jeder Euro wirkt gleich --- vom ersten bis zum dreißigsten. Das Experiment zeigt, dass es nicht stimmt.\\
\textbf{Ein Effekt} & Derselbe für Gesunde und Kranke, Arme und Reiche. Befund 4 widerspricht.\\
\textbf{Keine Nebenwirkung} & Zuzahlung ändert nur die Besuche --- nicht die Medikation, nicht den Zeitpunkt, nicht die Notaufnahme.\\
\textbf{Stabilität} & Der Effekt von 1975 gilt 2025, der amerikanische gilt in Deutschland.\\
\bottomrule
\end{tabular}}
\end{center}

\begin{center}
\begin{minipage}{0.88\textwidth}
\kursiv{\textbf{Die dritte Zeile ist die, an der Zuzahlungen in der Praxis scheitern}: Wer den Hausarztbesuch spart, kommt später in die Notaufnahme, und die kostet das Zehnfache. Ob die Gesamtrechnung aufgeht, weiß keine der Studien dieses Kapitels, weil keine das Gesamtsystem beobachtet. \textbf{Die Ökonometrie schätzt den Effekt auf die Größe, die sie misst} --- und die Entscheidung über die Zuzahlung braucht den Effekt auf alle.}
\end{minipage}
\end{center}
\end{warnung}

\section{Übungsaufgaben}

\begin{aufgabe}
\boxtitel{Aufgabe 14.1 --- Von der Elastizität zur Zahl}
Die Elastizität sei $-0{,}2$. Eine Kasse mit $5{,}0$ Besuchen je Versichertem erhöht die Zuzahlung von $10$ auf $15$ Euro. Was ist zu erwarten?

\loesung
\schritt{Schritt 1 --- Die Preisänderung}
\[ \frac{15-10}{10} = +50\,\% \quad\text{(Bogen: } \tfrac{5}{12{,}5} = 40\,\%\text{)} \]

\schritt{Schritt 2 --- Die Mengenänderung}
\[ -0{,}2\cdot40\,\% = -8\,\% \quad\Longrightarrow\quad 5{,}0\cdot0{,}92 = 4{,}6 \ \text{Besuche} \]

\schritt{Schritt 3 --- Die Annahmen in dieser Rechnung}
\begin{center}
{\small
\renewcommand{\arraystretch}{1.5}
\begin{tabular}{@{}P{2.8cm}P{5.8cm}@{}}
\toprule
\textbf{Konstante Elastizität} & gilt im Experiment nur grob; bei kleinen Zuzahlungen war sie höher.\\
\textbf{Übertragbarkeit} & vom Selbstbeteiligungssatz auf eine Pauschale, von 1975 auf heute.\\
\textbf{Keine Erwartung} & die Erhöhung wurde nicht angekündigt --- sonst Kapitel~13.\\
\bottomrule
\end{tabular}}
\end{center}

\begin{center}
\begin{minipage}{0.9\textwidth}
\textbf{Etwa $0{,}4$ Besuche weniger je Versichertem --- mit einem Vertrauensbereich, der aus der Übertragbarkeit stammt und den keine Statistik beziffert.}
\end{minipage}
\end{center}

\begin{center}
\begin{minipage}{0.88\textwidth}
\kursiv{\textbf{Das ist die Empfindlichkeitsanalyse aus Band 24 auf eine ökonometrische Schätzung angewandt}: Man variiere die Elastizität zwischen $-0{,}1$ und $-0{,}3$ und erhält $0{,}2$ bis $0{,}6$ Besuche. Diese Spanne ist der ehrliche Befund, und für die Budgetplanung der Kasse ist sie brauchbar --- für die Frage nach der Gesundheit der Versicherten nicht.}
\end{minipage}
\end{center}
\end{aufgabe}

\begin{aufgabe}
\boxtitel{Aufgabe 14.2 --- Die Kontrollgruppe wählen}
Für die Praxisgebühr stehen drei Kontrollgruppen zur Wahl. Bewerten Sie jede nach Kapitel~9.

\loesung
\schritt{Die drei Gruppen}
\begin{center}
{\small
\renewcommand{\arraystretch}{1.5}
\begin{tabular}{@{}P{2.6cm}P{3.0cm}P{2.8cm}@{}}
\toprule
\textbf{Gruppe} & \textbf{Für parallele Trends} & \textbf{Dagegen}\\
\midrule
Kinder & gleiche Kassen, gleiche Ärzte & eigene Vorsorgeprogramme; Eltern entscheiden\\
Privat\-versicherte & gleiche Altersstruktur & anderes Vergütungssystem; eigene Reformen\\
Befreite & gleiche Kassen & chronisch krank; Besuche kaum preisabhängig\\
\bottomrule
\end{tabular}}
\end{center}

\schritt{Die Prüfung}
\begin{center}
{\small
\renewcommand{\arraystretch}{1.5}
\begin{tabular}{@{}P{1.0cm}P{7.4cm}@{}}
\toprule
$1$ & Vorperioden $2000$--$2003$ für jede Gruppe: Welche entwickelte sich parallel?\\
$2$ & Ereignisstudie je Quartal: Setzt der Effekt zum Stichtag ein oder davor (Ankündigung, Kapitel~13)?\\
$3$ & Dreifache Differenz: Kinder in Reformkassen gegen Kinder anderswo --- falls es \enquote{anderswo} gibt.\\
\bottomrule
\end{tabular}}
\end{center}

\begin{center}
\begin{minipage}{0.9\textwidth}
\textbf{Keine Gruppe ist überzeugend. Der ehrliche Befund ist die Spanne über alle drei --- null bis zehn Prozent --- mit der Angabe, welche Annahme welchen Wert trägt.}
\end{minipage}
\end{center}

\begin{center}
\begin{minipage}{0.88\textwidth}
\kursiv{\textbf{Das ist kein Scheitern der Methode, sondern ihr Ergebnis}: Die Praxisgebühr traf alle Erwachsenen zugleich, und eine Maßnahme ohne Unbetroffene hat keine Kontrollgruppe. Wer das vorher weiß, kann die Einführung gestaffelt oder regional versetzt planen --- die Bundesländer hätten sich angeboten. Dass man es nicht tat, ist der Grund, warum die Frage nach zwanzig Jahren offen ist.}
\end{minipage}
\end{center}
\end{aufgabe}

\begin{aufgabe}
\boxtitel{Aufgabe 14.3 --- Eine Studie beurteilen}
Eine Studie findet aus Abrechnungsdaten einer Kasse, dass Versicherte im Hochzuzahlungstarif $30$ Prozent weniger Arztbesuche haben, kontrolliert für Alter, Geschlecht und dokumentierte Diagnosen. Was fragt man?

\loesung
\schritt{Die Fragen, nach Kapiteln}
\begin{center}
{\small
\renewcommand{\arraystretch}{1.5}
\begin{tabular}{@{}P{0.9cm}P{7.5cm}@{}}
\toprule
$1$ & Wer wählt den Tarif? Gesunde. Dann ist der Effekt überschätzt --- Richtung bekannt.\\
$1$ & Sind die Diagnosen eine Folge der Besuche? Ja --- dann verzerrt die Kontrolle zusätzlich.\\
$3$ & Gibt es eine Veränderung von außen --- Tarifreform, Stichtag?\\
$6$ & Wer ist in den Daten? Nur, wer in der Kasse blieb; Wechsler fehlen.\\
$7$ & Gibt es dieselben Versicherten vor und nach einem Tarifwechsel?\\
$13$ & Wurde der Tarif gewählt, um Besuche zu vermeiden --- oder umgekehrt?\\
\bottomrule
\end{tabular}}
\end{center}

\schritt{Das Urteil}
\begin{center}
\begin{minipage}{0.9\textwidth}
\textbf{Die $30$ Prozent sind eine Obergrenze. Der kausale Effekt liegt darunter, vermutlich weit --- und ohne eine Veränderung von außen lässt sich nicht sagen, wie weit.}
\end{minipage}
\end{center}

\begin{center}
\begin{minipage}{0.88\textwidth}
\kursiv{\textbf{Die Studie ist nicht wertlos}: Sie gibt eine Schranke, und sie gibt die Richtung der Verzerrung. Das ist mehr, als viele Studien bieten --- aber es ist weniger, als ihre Zusammenfassung behaupten wird. Der Leser dieses Bandes kann den Unterschied benennen, und das ist das Ziel.}
\end{minipage}
\end{center}
\end{aufgabe}

\begin{geschichte}
\boxtitelg{Vierhundert Millionen Dollar für eine Zahl}
Das Krankenversicherungsexperiment kostete nach heutigem Geld etwa $400$ Millionen Dollar --- und lieferte im Kern eine Elastizität.

\begin{center}
\begin{minipage}{0.9\textwidth}
\kursiv{\Person{Joseph Newhouse} schrieb 1993 die Zusammenfassung unter dem Titel \kursiv{Free for All?} \textbf{Er bezeichnete das Experiment später als das teuerste und zugleich billigste seiner Art}: teuer in absoluten Zahlen, billig gemessen an den Entscheidungen, die seither auf seinen Ergebnissen beruhen. Jede amerikanische Gesundheitsreform seit 1980 hat die Zahl $-0{,}2$ zitiert, meist von beiden Seiten.}
\end{minipage}
\end{center}

\schritt{Oregon}
\begin{center}
\begin{minipage}{0.9\textwidth}
\kursiv{2008 verloste der Bundesstaat Oregon Plätze in der Krankenversicherung für Arme --- weil das Geld nicht für alle reichte. \textbf{\Person{Amy Finkelstein} und Mitarbeiter nutzten die Lotterie als Instrument} und fanden: mehr Arztbesuche, weniger Depressionen, weniger finanzielle Not --- und nach zwei Jahren keine messbare Änderung bei Blutdruck, Cholesterin oder Blutzucker. Der Befund wird bis heute von beiden Seiten des Streits zitiert, wie der von 1982.}
\end{minipage}
\end{center}

\schritt{Und die Lehre}
\begin{center}
\begin{minipage}{0.9\textwidth}
\kursiv{\textbf{Beide Studien beantworten die Frage, die sie stellen konnten} --- die Wirkung auf Nutzung und auf wenige Gesundheitsmaße über wenige Jahre. Die Frage, die die Politik stellt --- lohnt sich Versicherung? ---, ist größer, und keine Methode dieses Bandes reicht an sie heran. Das ist kein Grund, die Methoden nicht zu benutzen. Es ist der Grund, sie richtig zu benennen.}
\end{minipage}
\end{center}
\end{geschichte}

\begin{forschung}
\boxtitelp{Die Zuzahlungsfrage heute}
\begin{center}
{\small
\renewcommand{\arraystretch}{1.45}
\begin{tabular}{@{}P{3.2cm}P{5.0cm}@{}}
\toprule
\textbf{Wertorientierte Zuzahlung} & Null für wirksame, hoch für unwirksame Leistungen --- Befund 2 umgehen.\\
\textbf{Verhaltens\-ökonomik} & Warum kleine Beträge große Wirkung haben --- Band 24, Kapitel~12.\\
\textbf{Verwaltungsdaten} & Ganze Länder als Panel, mit Reformen als Instrumenten.\\
\bottomrule
\end{tabular}}
\end{center}

\schritt{Zur zweiten Zeile}
\begin{center}
\begin{minipage}{0.88\textwidth}
\kursiv{Mehrere Studien seit 2015 finden, dass Versicherte bei Zuzahlungen auch lebenswichtige Medikamente absetzen --- \Person{Chandra}, \Person{Flack} und \Person{Obermeyer} zeigten 2021 für amerikanische Rentner einen messbaren Anstieg der Sterblichkeit nach einer kleinen Erhöhung der Zuzahlung. \textbf{Der Befund ist mit der Erwartungsnutzentheorie aus Band 24 schwer vereinbar} und gut mit der Neuen Erwartungstheorie: Der Verlust von zehn Dollar wiegt schwerer als das Risiko, das man nicht sieht. Für die Ökonometrie heißt das: Die Elastizität ist keine Konstante der Natur, sondern ein Ergebnis von Rahmung.}
\end{minipage}
\end{center}
\end{forschung}

\begin{merksatz}
\boxtitel{Was Sie jetzt können}
\begin{center}
\renewcommand{\arraystretch}{1.5}
\begin{tabular}{@{}P{3.0cm}P{5.6cm}@{}}
\toprule
\textbf{Die Tabelle} & zehn Verfahren, zehn Annahmen\\
\textbf{Das Experiment} & $-0{,}2$, und für wen es gilt\\
\textbf{Die Praxisgebühr} & die Kontrollgruppe entscheidet\\
\textbf{Die Berichtspflicht} & Annahme, Prüfung, Empfindlichkeit, Spanne\\
\textbf{Der Bruchfall} & Effekt auf das Gemessene, nicht auf alles\\
\bottomrule
\end{tabular}
\end{center}

\begin{center}
\begin{minipage}{0.88\textwidth}
\kursiv{Damit endet Teil~V. Das letzte Kapitel zieht Bilanz --- und fragt, was die Ökonometrie den Nachbarfächern der Reihe gegeben hat, und was sie von ihnen bekommen hat.}
\end{minipage}
\end{center}
\end{merksatz}
